\subsubsection{The selected IMT maps}

The state dimension is $s=19$ and the step width is $t=128$ throughout this
appendix. Table~\ref{tab:imt-maps} specifies both maps in
\eqref{eq:structured-inner}. Bit $i$ of each hexadecimal word is coordinate
$i$, starting with the least significant bit. Write $M:=2^{19}-1$.
The nonzero image weights of $A$ and their multiplicities are
\[
 (d_i)_{i=1}^5=(48,56,64,72,80),\qquad
 (c_i)_{i=1}^5=(5166,110288,293455,110128,5250).
\]
Every coordinate of $A$ is a nonzero linear functional. The $128$ columns
of $C$ are distinct nonzero vectors of weight five and span $\F_2^{19}$.
The artifact checks these properties and both spectra by exact enumeration.

\begin{table}[t]
\centering
\caption{Fixed IMT expansion and feedback maps. The words specify
$A(e_j)$ and $C^\top(e_j)$ for the standard state basis.}
\label{tab:imt-maps}
\scriptsize
\setlength{\tabcolsep}{3pt}
\begin{tabular}{@{}rll@{}}
\toprule
$j$ & $A(e_j)$ & $C^\top(e_j)$ \\
\midrule
0 & \texttt{55555555555555555555555555555555} & \texttt{800003d868100423210234004341a004} \\
1 & \texttt{33333333333333333333333333333333} & \texttt{010d0102568441441a6212120884488c} \\
2 & \texttt{0f0f0f0f0f0f0f0f0f0f0f0f0f0f0f0f} & \texttt{40482008104050000312449182000a43} \\
3 & \texttt{ff00ff00ff00ff00ff00ff00ff00ff00} & \texttt{c60682141218101b9a382140d1011330} \\
4 & \texttt{0000ffff0000ffff0000ffff0000ffff} & \texttt{02800a0201093ca2cc121a20200c0408} \\
5 & \texttt{ffffffff00000000ffffffff00000000} & \texttt{81212844806ef00068055044102ee101} \\
6 & \texttt{ffffffffffffffff0000000000000000} & \texttt{1e200091046309f08018083800400061} \\
7 & \texttt{ca9ff6a3c56f06ac0653c590f65cca60} & \texttt{91750062141381019100a60c30a810aa} \\
8 & \texttt{8e8e718e2bd4d4d4b2b24db217e8e8e8} & \texttt{01043040c08a882804068c12246500d0} \\
9 & \texttt{b2b2b2b218e7e71842bdbd42e8e8e8e8} & \texttt{2089148c8300089c8504044184424100} \\
10 & \texttt{a00a05af366393c6396c633650fa0aa0} & \texttt{044340810132201506e1012001322008} \\
11 & \texttt{9350635fa063506c5f9c506c93509ca0} & \texttt{42001d281e80022101ec18224403cc12} \\
12 & \texttt{12e2747b121d7484de2e4748212eb848} & \texttt{2834da0440600280021100010b9a9780} \\
13 & \texttt{4411e14b11bb4b1e2d87772287d22288} & \texttt{2118a12209240f06640041cacb0884c5} \\
14 & \texttt{59953003cffca66a0cc0655665560cc0} & \texttt{80122401c000804c00c9821294803a04} \\
15 & \texttt{7447b88b74b8b874b7847b4848848448} & \texttt{6880cc1f224046004080818c3814003b} \\
16 & \texttt{11ddd21e447787b42d1e112278b44488} & \texttt{1cf284802d0ca24a000820040c904700} \\
17 & \texttt{72e48214e48deb82d8b1d7be4ed8be28} & \texttt{7c405360a885745070a0c0a122702002} \\
18 & \texttt{4e824e82e428e4284e824e82e428e428} & \texttt{128a40312091018038456bcd40011874} \\
\bottomrule
\end{tabular}
\end{table}

\subsubsection{Finite IMT transfers}
\label{app:imt-finite-transfers}

We represent a weighted state measure by seven nonnegative coordinates
$(Z,D,S_1,\ldots,S_5)$. The $Z$ coordinate is mass at zero. The $D$
coordinate bounds arbitrary mass on nonzero states. Coordinate $S_i$
dominates $S_i$ times the uniform probability measure on
$\{q:\wt(Aq)=d_i\}$. Terminal mass is bounded by the sum of coordinates.
Rows of a transfer matrix index the entering coordinate.

First fix an input weight $j$. Let $n_j:=\binom{128}{j}$, let $k_j$ be the
number of weight-$j$ vectors in $\ker C$, and put $\nu_j:=1-k_j/n_j$.
For $z\in(0,1)$, the emitted moment at image weight $d$ is
\[
 m_{d,j}(z):=\frac1{n_j}\sum_h
 \binom dh\binom{128-d}{j-h}z^{d+j-2h}.
\]
Set $m_{D,j}:=\max_i m_{d_i,j}$. For $q\ne0$, define the lazy
cancellation moment
\[
 L_j(q;z):=\frac1{n_j}\sum_{\substack{\wt(X)=j\\CX=q}}
 z^{\wt(X+Aq)}.
\]
Choose upper bounds $\ell_{D,j}\ge\max_{q\ne0}L_j(q;z)$ and
$\ell_{i,j}\ge c_i^{-1}\sum_{\wt(Aq)=d_i}L_j(q;z)$.
Let $f_j$ bound the size of every nonzero weight-$j$ syndrome fiber. Valid
choices are the minima of the following bounds:
\begin{align*}
 \ell_{D,j}&:=
 \min\left\{m_{D,j},\nu_j,
 \frac{f_j}{n_j}z^{\min_i|d_i-j|}\right\},\\
 \ell_{i,j}&:=
 \min\left\{m_{d_i,j},
 \frac{\min(n_j-k_j,c_i f_j)}{c_i n_j}z^{|d_i-j|}\right\}.
\end{align*}
For $j=1,2$ the verifier also uses the exact cancellation sums, enumerating
the input supports and grouping them by $CX$ and $\wt(ACX)$. No relation
between $A$ and $C^\top$ is assumed.

The integer fiber bounds require only the spectrum of $C^\top$. Let $b_w$
count its nonzero words of weight $w$, and let $\mathcal K_j(w)$ be the
binary Krawtchouk polynomial. Fourier inversion and Parseval give
\[
 k_j=2^{-19}\left(n_j+\sum_w b_w\mathcal K_j(w)\right),\qquad
 V_j=2^{-19}\left(n_j^2+\sum_w b_w\mathcal K_j(w)^2\right),
\]
where $V_j$ is the sum of squares of all syndrome fiber sizes. Put
$a_j:=n_j-k_j$. A nonzero fiber is at most each of
\begin{gather*}
 a_j,\qquad
 \left\lfloor2^{-19}\left(n_j+\sum_w b_w|\mathcal K_j(w)|\right)\right\rfloor,\\
 \left\lfloor\frac{a_j+\lceil\sqrt{(M-1)(M(V_j-k_j^2)-a_j^2)}\rceil}{M}\right\rfloor.
\end{gather*}
The last bound follows by minimizing the other $M-1$ squares for a fixed
maximum fiber. There is also a constant-weight packing bound. If $d_C$ is
the kernel distance, set $r=\lfloor(d_C-1)/2\rfloor$ and $h=\min(j,128-j)$.
The bound is $1$ for $h<r$ and
$\lfloor\binom{128}{h-r}/\binom hr\rfloor$ otherwise: two words in one
fiber cannot contain the same $(h-r)$-subset after complementing when needed.
The verifier takes the minimum of these integer upper bounds.

Define the fixed-weight transfer $T_j(z)$ by the following nonzero entries:
\begin{align*}
 (T_j)_{ZZ}&=(1-\nu_j)z^j,& (T_j)_{ZD}&=\nu_j z^j,\\
 (T_j)_{DZ}&=\tfrac12\ell_{D,j}+\tfrac1{2M}\min(m_{D,j},\nu_j),&
 (T_j)_{DD}&=\tfrac12m_{D,j},\\
 (T_j)_{DS_i}&=\tfrac{c_i}{2M}m_{D,j},&
 (T_j)_{S_iZ}&=\tfrac12\ell_{i,j}+\tfrac1{2M}\min(m_{d_i,j},\nu_j),\\
 (T_j)_{S_iS_h}&=\tfrac{c_h}{2M}m_{d_i,j}.&&
\end{align*}
Additionally put mass $m_{d_i,j}/2$ in entry $S_iD$ when $\nu_j>0$, or
add it to entry $S_iS_i$ when $\nu_j=0$.

These entries dominate the weighted state measure. The zero row is exact.
The lazy branch pays its cancellation bound and places the remaining mass
in $D$; if all syndromes vanish, it preserves the entering shell.
On the refresh branch, every specified target has probability at most $1/M$.
Termination also requires a nonzero syndrome, giving the minimum in the
zero column. Keeping the full surviving-mass upper bounds deliberately
overcounts cancellation; it never subtracts an unproved lower bound.
For iid Bernoulli-$\beta$ input, the occupation transfer is
\[
 T_{\rm occ}(\beta,z):=\sum_{j=0}^{128}
 \binom{128}{j}\beta^j(1-\beta)^{128-j}T_j(z).
\]

A second envelope bounds every target by Fourier inversion of a tilted
input law. Set
\[
 g_0:=1-\beta+\beta z,\quad g_1:=\beta+(1-\beta)z,\quad
 F_d:=g_0^{128-d}g_1^d,
\]
\[
 r_0:=|1-2\beta z/g_0|,\qquad
 r_1:=|1-2(1-\beta)z/g_1|,
\]
and define
\[
 H_d:=2^{-19}\left(1+\sum_w b_w
 \max_{\max(0,d+w-128)\le h\le\min(d,w)}r_0^{w-h}r_1^h\right).
\]
The maximum is attained at an endpoint. Tilting by $z^{\wt(X+Aq)}$ keeps
the input coordinates independent. Fourier inversion bounds each syndrome
point probability by $H_d$. Thus every target, including zero, receives
weighted mass at most $F_d(H_d/2+1/(2M))$. Define $T_{\rm F}$ by placing
this value in column $Z$ and $c_h$ times this value in column $S_h$ for an
entering shell $S_i$ with $d=d_i$. For an entering $D$, maximize over $d_i$.
Its zero row is the exact binomial mixture of the zero rows of $T_j$.

Finally the state-independent scalar envelope is
\[
 T_{\rm sc}(\beta,z):=\max_{d\in\{0,48,56,64,72,80\}}F_d.
\]
At $\beta=1/2$ it equals $((1+z)/2)^{128}$ exactly. For either matrix
envelope, induction gives
\begin{equation}
 \E[z^{\wt(Y)}]\le e_Z^\top T(\beta,z)^{N/128}\1.
 \label{eq:imt-moment}
\end{equation}
The scalar bound gives $T_{\rm sc}^{N/128}$. A witness selects one complete
envelope; the proof never takes an entrywise minimum of representations.

\subsubsection{Positive occupation}

Fix $\alpha\ge10^{-4}$ and mean active-row weight $x\in\mathcal W$.
All arrays in a fixed profile have $N\alpha x$ ones. Apply route domination
to the moment restricted to that weight slice, then change its iid reference
probability from $\alpha x$ to $py$. The likelihood cost is
$\exp(ND_{\rm KL}(\alpha x\Vert py))$. The chain rule bounds its exponent by
$D_{\rm KL}(\alpha\Vert p)+\alpha D_{\rm KL}(x\Vert y)$.
Consequently the bad-message exponent per output bit is at most
\begin{equation}
 \alpha\widehat a_{\rm BA}(x)+D_{\rm KL}(\alpha\Vert p)
 +\alpha D_{\rm KL}(x\Vert y)+\frac{\ln\lambda}{128}-0.11\ln z,
 \label{eq:imt-dense-exponent}
\end{equation}
where $\lambda$ is a scalar bound or a positive Collatz bound
$T(py,z)\boldsymbol w\le\lambda\boldsymbol w$.

The verifier covers $[10^{-4},1]\times\mathcal W$ by $1{,}023$ rational
boxes within the $39$ outer segments. It fixes a witness on each box and
checks four vertices. The exponent is convex in $(\alpha,\alpha x)$ on
each affine segment. All bounds pass outward arithmetic at $256$ bits and
replay at $512$ bits, with maximum less than $-4.10335\cdot10^{-7}$.
The witnesses use $433$ occupation transfers, $307$ Fourier transfers, and
$283$ scalar transfers. The values $p,y$ and the positive vectors are rational;
the weight tilt is specified by a rational $\log(-\log z)$ and enclosed
outward. Exact rational geometry checks containment, disjoint interiors,
and total area. No numerical eigenvalue is trusted.

There are finitely many witnesses, so their moment prefactors are bounded.
Counting and route-domination costs are $\exp(o(N))$ uniformly in this range.
They include active-position choices, row-weight tuples, and the selected
spectrum's polynomial factors. The summed first moment is therefore $o(1)$.

\subsubsection{Uniform sparse occupation}

For sparse profiles, the factor $(L+1)^b$ would be too costly. We instead
compare the counting measure on each permuted active outer row with fair
bits. The original, coarser $31$-segment outer majorant already proves
\[
 \sup_{x\in\mathcal W}\bigl(\widehat a_{\rm old}(x)-h(x)+\ln(2)/2\bigr)
 <\ell,\qquad \ell:=0.01281.
\]
Hence, uniformly over all selected outers and weights, for some
$\epsilon_b=O(\log(b)/b)$,
\begin{equation}
 \frac{A_{\mathcal O_b}(w)}{\binom bw}
 \le \exp\{b(-\ln(2)/2+\ell+\epsilon_b)\}.
 \label{eq:imt-row-likelihood}
\end{equation}
Relative to a fair $b$-bit row, the counting measure thus costs
$\exp\{b(\ln(2)/2+\ell+\epsilon_b)\}$. This comparison is made after
summing messages; it does not assert that a fixed constituent word is uniform.

Fix $Q$ active positions and put $\alpha=Q/L\le10^{-4}$.
Under the fair-row reference, each region has $Q$ uniform marked positions
carrying independent fair bits. Sample iid marks with probability
$p=(8/5)\alpha$ and condition their count to equal $Q$. Without that
conditioning, all input bits are iid with $\beta=p/2=(4/5)\alpha$.
The reference regions are independent. Since $Q$ is a mode of
$\operatorname{Bin}(L,\alpha)$, Chebyshev's inequality gives modal mass
at least $1/(8\sqrt Q)$. The exact binomial likelihood ratio yields
\[
 \Prob[\operatorname{Bin}(L,p)=Q]
 \ge \frac{e^{-LD_{\rm KL}(\alpha\Vert p)}}{8\sqrt Q}.
\]
Indeed, probability at least $3/4$ lies in at most $4\sqrt Q+1\le5\sqrt Q$
integer locations about the mean, so a mode has mass at least $3/(20\sqrt Q)$.

Put $z=1-(8/5)\alpha$. An exact rational polynomial certificate proves
\begin{equation}
 T_{\rm occ}(\beta,z)\boldsymbol w(\alpha)
 \le (1-96\alpha)\boldsymbol w(\alpha),
 \qquad 0<\alpha\le10^{-4}.
 \label{eq:imt-sparse-collatz}
\end{equation}
Here $w_Z=1$, and $w_i=2^{-10}(1+h_i\alpha)$ for the six live coordinates,
with
\begin{align*}
 h_D&=900,& h_{S_1}&=\frac{116184452864}{2257055535},&
 h_{S_2}&=\frac{92653723520}{3613910291},\\
 h_{S_3}&=-\frac{26853143552}{769273207925},&
 h_{S_4}&=-\frac{461216096896}{18043337105},&
 h_{S_5}&=-\frac{351546952448}{6881266875}.
\end{align*}
The shell-weighted average of the corrections is zero. Every coordinate is
at least $1/2048$. In the variable $x=10000\alpha$, the verifier constructs
degree-$257$ rational polynomial upper bounds for the seven residuals.
It certifies every maximum replacement over $[0,1]$. Each residual vanishes
at zero. After removing that factor, its constant coefficient plus the sum
of its positive remaining coefficients is strictly negative. This proves
the whole interval, including arbitrarily small positive $\alpha$.

Combining \eqref{eq:imt-row-likelihood}, the conditioning bound,
\eqref{eq:imt-sparse-collatz}, and Markov's inequality gives
\begin{align}
 \E[Z_{\lfloor0.11N\rfloor,Q}\mid\mathcal G_b]
 \le{}&2048\binom LQ e^{Qb(\ln(2)/2+\ell+\epsilon_b)}(8\sqrt Q)^b
 \nonumber\\*
 &\cdot e^{ND_{\rm KL}(\alpha\Vert p)}
 (1-96\alpha)^{N/128}z^{-0.11N}.
 \label{eq:imt-uniform-sparse}
\end{align}
Use $\ln\binom LQ\le Q\ln L+Q$, $b\ge c\log_2N$ with $c=39/4$, and
$\ln(1-u)\le-u$. The uniform coefficient per active-row bit is at most
\[
 C_*:=\frac{\ln2}{2}+\ln(5/8)
 +\frac{3/5+0.11(8/5)}{1-0.00016}-\frac{96}{128}
 +\ell+\frac{\ln2}{c}<-0.01340384.
\]
The remaining terms are
$\epsilon_b+1/b+\ln(8\sqrt Q)/Q+\ln(2048)/(Qb)$.
For $Q\ge4096$, the third is below $0.002$ and decreases with $Q$.
For sufficiently large $b$, the other terms together are below $0.004$
uniformly. Thus
\begin{equation}
 \E[Z_{\lfloor0.11N\rfloor,Q}\mid\mathcal G_b]
 \le e^{-0.006Qb},\qquad 4096\le Q\le10^{-4}L.
 \label{eq:imt-sparse-final}
\end{equation}
This explicit cutoff avoids any assumption about how fast $Q$ grows.

\subsubsection{Fixed occupation and the continuum limit}

It remains to control finitely many $Q<4096$. We derive a limit for each
fixed $Q$, allowing constants to depend on $Q$. Define
\[
 r:=1/M,\qquad p_0:=2^{18}/M,\qquad
 P_0:=\begin{pmatrix}0&1\\r&1-r\end{pmatrix}.
\]
During empty epochs a live state has transition $P=(I+\Pi)/2$, where
$\Pi$ replaces it by a uniform nonzero state. Write $f(q):=\wt(Aq)$.
Its stationary mean is $\mu=128p_0$. For any initial live state,
\[
 \E[f(q_j)\mid q_i]=\mu+2^{-(j-i)}(f(q_i)-\mu),\qquad j\ge i.
\]
For $S_g:=\sum_{i=0}^{g-1}f(q_i)$ this implies
\[
 |\E S_g-\mu g|\le2t,\qquad
 \operatorname{Var}(S_g)\le3t^2g,\qquad
 \E|S_g-\mu g|\le t\sqrt{3g}+2t.
\]
The variance bound follows from
$\operatorname{Cov}(f(q_i),f(q_j))=2^{-(j-i)}\operatorname{Var}(f(q_i))$.
It does not require a stationary initial state.

Fix $\theta>0$, and set $z_L=e^{-\theta/L}$ and $\gamma=p_0\theta$.
Let $W_0(z_L)$ be the actual full-state weighted empty-step kernel.
The Lipschitz bound for $e^{-x}$ and
$P^g=2^{-g}I+(1-2^{-g})\Pi$ give, for every live $q,q'$,
\begin{equation}
 \left|W_0(z_L)^g(q,q')-\frac{e^{-\theta\mu g/L}}M\right|
 \le\frac{\theta t(\sqrt{3g}+2)}L+2^{-g}.
 \label{eq:imt-empty-limit}
\end{equation}
The zero state stays zero. Define the lift $J$ by $J(q,0)=\1_{q=0}$,
$J(q,1)=\1_{q\ne0}$, and the projection $R$ by
$R(0,q)=\1_{q=0}$, $R(1,q)=\1_{q\ne0}/M$. Thus $RJ=I_2$.
Put $D_\gamma(u):=\operatorname{diag}(1,e^{-\gamma u})$.
For $\lceil3\log_2L\rceil\le g\le L/t$, equation
\eqref{eq:imt-empty-limit} approximates $W_0^g$ by
$JD_\gamma(tg/L)R$ with entry error $O(L^{-1/2})$.

Let $W_1(z_L)$ be the kernel for one input bit at a uniformly chosen local
position. Its difference from $W_1(1)$ is $O(1/L)$ in row norm, since the
emitted weight is at most $t$. Invertibility of the mixer and the nonzero
columns of $C$ give the exact identity
\[
 RW_1(1)J=P_0.
\]
For a fixed $a\le Q$ uniformly placed impulses in one region, an epoch
collision has probability $O_Q(1/L)$. A boundary or inter-impulse gap shorter
than $\lceil3\log_2L\rceil$ has probability $O_Q(\log(L)/L)$.
On the complement, replace each empty gap using
\eqref{eq:imt-empty-limit} and each impulse by $W_1(1)$. The finite products
telescope with bounded row norms. The omitted impulse durations total
$at/L=O_Q(1/L)$. Riemann sums for the ordered uniform sites then give
\begin{equation}
 K_a(\theta):=a!\int_{\substack{u_i\ge0\\\sum_{i=0}^a u_i=1}}
 D_\gamma(u_0)P_0D_\gamma(u_1)\cdots P_0D_\gamma(u_a)\,d\boldsymbol u.
 \label{eq:imt-continuum}
\end{equation}
For $a=0$, $K_0=D_\gamma(1)$. The simplex integral uses $a$ independent
coordinates. The actual region kernel $F_{a,L}(z_L)$ differs from $JK_aR$
by $O_{Q,M,\theta}(L^{-1/2}+\log(L)/L)$ in every entry, uniformly over
initial and final states. Bad placement events cost at most their probability
because weighted actual kernels have row sums at most one.

The additive comparison becomes multiplicative. For $a=0$, both zero/live
cross entries vanish exactly. For $a=1$, zero-to-zero vanishes exactly.
All other coarse entries are positive, and for $a\ge2$ every entry is
positive. Since $Q,M,\theta$ are fixed, there exists
$\eta_L=O_{Q,M,\theta}(L^{-1/2}+\log(L)/L)$ such that
\begin{equation}
 F_{a,L}(z_L)\le(1+\eta_L)JK_a(\theta)R\quad(0\le a\le Q).
 \label{eq:imt-componentwise-limit}
\end{equation}
This inequality averages over the within-region placements. It is uniform
over the choices of which active rows contribute a bit and over boundary
states. Multiplication over $b=O(\log L)$ regions costs $\exp(o(1))$,
also when coefficient extraction enforces the row weights.

We may enlarge $P_0$ to
$P_+:=\left(\begin{smallmatrix}0&1\\r&1\end{smallmatrix}\right)$
in \eqref{eq:imt-continuum}; write $K_a^+$ for the resulting upper matrices.
They are nonnegative, but not stochastic.

\paragraph{One and two active rows.}
Use \eqref{eq:imt-row-likelihood} and fair reference rows. The continuum
region matrix is $2^{-Q}T_Q$, with
$T_Q:=\sum_{a=0}^Q\binom Qa K_a^+$. Set $\gamma=Q\ln2$ and
$\theta=\gamma/p_0$. With
\[
 a_0=e^{-\gamma},\quad f=\frac{1-a_0}{\gamma},\quad
 g=\frac{2(\gamma-1+a_0)}{\gamma^2},\quad
 e=\frac{2(1-(1+\gamma)a_0)}{\gamma^2},
\]
the matrices are
\[
 K_0^+=\operatorname{diag}(1,a_0),\quad
 K_1^+=\begin{pmatrix}0&f\\rf&a_0\end{pmatrix},\quad
 K_2^+=\begin{pmatrix}rg&e\\re&re+a_0\end{pmatrix}.
\]
The exponent per outer bit, including active-position choices, is
\[
 \ln\rho(T_Q)+0.11\theta-Q\ln(2)/2+Q\ell+Q\ln(2)/c.
\]
Outward checks at $256$ bits, replayed at $512$ bits, bound this expression
divided by $Q$ by $-0.10918$ for $Q=1$ and $-0.10914$ for $Q=2$.
Finite positive-matrix prefactors and the uniform outer remainder are
$\exp(o_Q(b))$, so both first moments tend to zero.

\paragraph{A weight-coupled bound for fixed $Q\ge3$.}
For positive row fugacities $u_j$, define
\[
 T_Q(\boldsymbol u,\theta):=\sum_{\boldsymbol e\in\{0,1\}^Q}
 \left(\prod_j u_j^{e_j}\right)K_{\sum_j e_j}^+(\theta).
\]
Nonnegative coefficient extraction and \eqref{eq:imt-componentwise-limit}
bound the moment for row weights $w_j$ by
\[
 \left(\prod_j\binom b{w_j}^{-1}u_j^{-w_j}\right)
 e_0^\top T_Q(\boldsymbol u,\theta)^b\1_2\,\exp(o_Q(b)).
\]
Set $\sigma=127/250$, $\tau=133/125$, $v_1=3/1600$, and
$\theta=Q\tau/p_0$. Define
\begin{align*}
 G_Q&:=\sup_{0\le y\le1}\{-\tau y+(1+1/Q)\ln(1-y+y/\sigma)\},\\
 m_v(u)&:=\max\{1+u\sigma v_1,ur/v_1+\sigma(1+u)\}.
\end{align*}
For the weighted row norm with $v=(1,v_1)^\top$,
\begin{equation}
 \|T_Q(\boldsymbol u,Q\tau/p_0)\|_v
 \le e^{QG_Q}\prod_jm_v(u_j),\qquad G_Q\le G_3.
 \label{eq:imt-product-norm}
\end{equation}
For completeness, insert all $Q$ potential impulse sites, using $I$ for
an unused site and $P_+$ for a used site. If a fixed class path has $l$ live
spacings, their total duration $Y$ has distribution
$\operatorname{Beta}(l,Q+1-l)$, with deterministic endpoint cases. The identity
\[
 \E[(1-Y+Y/\sigma)^{-(Q+1)}]=\sigma^l
\]
follows by the substitution $u=y/[\sigma+(1-\sigma)y]$ in the beta integral.
By the definition of $G_Q$, the path's integrated weight is at most
$e^{QG_Q}\sigma^l$. Summing paths yields the average over orders $\pi$ of
\[
 D_\sigma(I+u_{\pi(1)}P_+)D_\sigma\cdots
 (I+u_{\pi(Q)}P_+)D_\sigma,
\]
times $e^{QG_Q}$, where $D_\sigma=\operatorname{diag}(1,\sigma)$.
Submultiplicativity, $\|D_\sigma\|_v\le1$, and
$\|(I+uP_+)D_\sigma\|_v=m_v(u)$ prove \eqref{eq:imt-product-norm}.
This reasoning requires nonnegative coefficients, not stochasticity of $P_+$.

The old $31$-segment majorant suffices here. For each segment choose one
rational $u>0$. Convexity in $x$ reduces the following inequality to its
two endpoints:
\[
 \widehat a_{\rm old}(x)-h(x)-x\ln u+\ln m_v(u)
 +G_3+0.11\tau/p_0+\ln(2)/c<-0.0008679786.
\]
All $62$ outward checks pass. The row-weight tuples and finite coefficient
remainders contribute $o_Q(b)$. Therefore, for each fixed $Q\ge3$,
\[
 \E[Z_{\lfloor0.11N\rfloor,Q}\mid\mathcal G_b]
 \le\exp\{-0.0008679Qb+o_Q(b)\}.
\]

\subsubsection{Completion and certificate boundary}

\begin{proof}[Theorem~\ref{thm:structured-spin-scalable}]
Condition on any outer satisfying $\mathcal G_b$. The fixed-occupation
arguments give a vanishing sum over the finite set $1\le Q<4096$.
Equation~\eqref{eq:imt-sparse-final} gives a uniform geometric sum for
$4096\le Q\le10^{-4}L$. The positive-occupation certificate bounds the
remaining sum by $L\exp(-\eta N+o(N))$ for some $\eta>0$.
Thus $\sum_{Q=1}^L\E[Z_{\lfloor0.11N\rfloor,Q}\mid\mathcal G_b]=o(1)$.
The explicit cutoff makes this a summable union, not merely a statement
about each occupation sequence. Markov's inequality and
$\Prob[\mathcal G_b^c]=o(1)$ prove the distance claim.

Rate follows from injectivity. Every outer row costs $O(b)$ in either circuit
direction: it uses fixed Golay circuits, two permutations, and two accumulators.
The route costs $O(N)$, and the fixed-size binary IMT steps cost $O(1)$ each.
The displayed adjoint gives the same bound for transposed encoding.
\end{proof}

The imported outer event and its coarse majorant retain the $31$-entry
manifest in the
\artifactlink{https://github.com/ladnir/permute_conv/blob/main/workstreams/paper_architecture/certificates/single_sampled_ba_rm2sub/README.md}{outer certificate snapshot}.
The IMT evidence has separate source bindings, including the refined
outer and dense replay.
The \artifactlink{https://github.com/ladnir/permute_conv/blob/main/workstreams/inner_design/imt_asymptotic/README.md}{IMT reproduction guide}
identifies the source bindings and commands. The numerical certificates check
the finite spectra, interval covers, Collatz inequalities, exact sparse
polynomials, and fixed-occupation constants. The analytic steps are the
measure comparisons, transfer induction, continuum limit, and final union
proved above. Higher-precision replay is not an independent proof of those
steps. The finite BCH/IMT results use the same transfer interface, with
their separate finite-length evidence in
Appendix~\ref{app:finite-bch-proof}.
